% =====================================================================
%  UML-Sequenzdiagramm (UML 2.5.1, Kap. 17): „Aktivität speichern ->
%  Dashboard aktualisiert sich“. Abgelesen aus dem eingefrorenen Code:
%  AktivitaetDialog._speichern, AktivitaetService.speichern,
%  AktivitaetRepository.speichern, Seite._geaendert,
%  DashboardSeite.aktualisieren, AuswertungService.wochenbewertung/hinweise.
%  Export: docs/diagramme/export/sequenz.pdf
% =====================================================================
\documentclass[tikz,border=4pt]{standalone}
\usepackage{iu-diagramm}

\tikzset{
  kopf/.style={draw, fill=iuFlaeche, align=center, minimum height=0.85cm, text width=1.55cm,
               inner sep=2pt, font=\fontsize{7.3pt}{8.4pt}\selectfont, anchor=south},
  leben/.style={dashed, draw=iuGrau},
  akt/.style={draw, fill=white, minimum width=0.18cm, inner sep=0pt},
  sync/.style={-{Triangle[length=2mm,width=1.6mm]}},
  async/.style={-{Straight Barb[length=2mm]}},
  antwort/.style={dashed, -{Straight Barb[length=2mm]}},
  lab/.style={font=\fontsize{7pt}{8pt}\selectfont, inner sep=1pt, fill=white, fill opacity=0.85,
              text opacity=1},
  rahmen/.style={draw, inner sep=0pt},
  waechter/.style={font=\fontsize{7pt}{8pt}\selectfont, inner sep=1.5pt},
}
% Lebenslinien: x-Positionen. Repository und Datenbank stehen rechts von
% Zustand, Dashboard und AuswertungService, damit lange Nachrichten keine
% aktiven Ausführungsbalken kreuzen.
\def\xN{0}\def\xD{1.95}\def\xS{3.95}\def\xV{5.85}\def\xZ{7.75}\def\xH{9.55}
\def\xA{11.3}\def\xR{13.05}\def\xB{14.85}\def\xW{16.6}\def\xG{18.3}
% Nachricht von x#1 nach x#2 auf Höhe #3 mit Beschriftung #4 (Stil #5),
% Beschriftung mittig über dem Pfeil
\newcommand{\nachricht}[5][sync]{%
  \draw[#1] (#2,#4) -- (#3,#4) node[lab, above, midway] {#5};}
% Aufruf von links nach rechts: Beschriftung über dem Pfeil, beginnend am
% Balken des Senders (der Balken des Empfängers beginnt erst auf Höhe des Pfeils)
\newcommand{\aufrufR}[5][sync]{%
  \draw[#1] (#2,#4) -- (#3,#4);
  \node[lab, anchor=south west] at (#2+0.04,#4) {#5};}
% Rückgabe: Beschriftung unter dem Pfeil, beginnend hinter der Pfeilspitze am
% Balken des Empfängers (der Balken des Senders endet auf Höhe des Pfeils)
\newcommand{\rueckgabeL}[4]{%
  \draw[antwort] (#1,#3) -- (#2,#3);
  \node[lab, anchor=north west] at (#2+0.25,#3-0.04) {#4};}
% Nachricht an sich selbst bei x#1, Höhe #2, Beschriftung #3
\newcommand{\selbst}[3]{%
  \draw[sync] (#1+0.09,#2) -- ++(0.45,0) -- ++(0,-0.32) -- ++(-0.36,0);
  \node[lab, anchor=west] at (#1+0.6,#2-0.16) {#3};}
\newcommand{\aktiv}[3]{\draw[fill=white] (#1-0.09,#2) rectangle (#1+0.09,#3);}

\begin{document}
\begin{tikzpicture}

% ---------------------------------------------------------------------
%  Köpfe und Lebenslinien
% ---------------------------------------------------------------------
\akteur{nutzer}{\xN,0.45}{}
\node[font=\fontsize{7.3pt}{8.4pt}\selectfont, anchor=south] at (\xN,1.05) {Nutzer:in};
\node[kopf] at (\xD,0) {:Aktivitaet\-Dialog};
\node[kopf] at (\xS,0) {:Aktivitaet\-Service};
\node[kopf] at (\xV,0) {«module»\\validierung};
\node[kopf] at (\xR,0) {:Aktivitaet\-Repository};
\node[kopf] at (\xB,0) {:Datenbank};
\node[kopf] at (\xZ,0) {:Zustand};
\node[kopf] at (\xH,0) {:Dashboard\-Seite};
\node[kopf] at (\xA,0) {:Auswertung\-Service};
\node[kopf] at (\xW,0) {«module»\\bewertung};
\node[kopf] at (\xG,0) {«module»\\regeln};
\foreach \x in {\xD,\xS,\xV,\xR,\xB,\xZ,\xH,\xA,\xW,\xG} { \draw[leben] (\x,0) -- (\x,-16.6); }
\draw[leben] (\xN,-0.1) -- (\xN,-16.6);

% ---------------------------------------------------------------------
%  Ausführungsbalken
% ---------------------------------------------------------------------
\aktiv{\xD}{-0.55}{-16.1}
\aktiv{\xS}{-1.65}{-5.45}
\aktiv{\xV}{-1.95}{-2.3}
\aktiv{\xR}{-3.4}{-4.95}
\aktiv{\xB}{-3.95}{-4.45}
\aktiv{\xZ}{-6.0}{-13.3}
\aktiv{\xH}{-6.55}{-13.15}
\aktiv{\xA}{-7.75}{-10.35}
\aktiv{\xR}{-8.3}{-8.8}
\aktiv{\xW}{-9.5}{-10.0}
\aktiv{\xA}{-11.1}{-12.45}
\aktiv{\xG}{-11.45}{-11.95}
\aktiv{\xS}{-14.35}{-14.6}

% ---------------------------------------------------------------------
%  Nachrichten
% ---------------------------------------------------------------------
\nachricht[async]{\xN}{\xD-0.09}{-0.55}{Speichern}
\selbst{\xD}{-0.8}{\_pruefen() : keine Fehler}
\nachricht{\xD+0.09}{\xS-0.09}{-1.65}{speichern(a)}
\aufrufR{\xS+0.09}{\xV-0.09}{-1.95}{pruefe\_aktivitaet(a, \dots)}
\rueckgabeL{\xV-0.09}{\xS+0.09}{-2.3}{fehler : list[Fehler]}

% alt-Fragment
\draw[rahmen] (\xD-0.5,-2.75) rectangle (\xG+0.75,-16.25);
\draw (\xD-0.5,-3.05) -- (\xD+0.15,-3.05) -- (\xD+0.3,-2.9) -- (\xD+0.3,-2.75);
\node[waechter, anchor=north west] at (\xD-0.48,-2.77) {\textbf{alt}};
\node[waechter, anchor=west] at (\xD+0.45,-2.98) {[fehler ist leer]};
\draw[dashed] (\xD-0.5,-13.85) -- (\xG+0.75,-13.85);
\node[waechter, anchor=west] at (\xD+0.45,-14.07) {[sonst]};

\nachricht{\xS+0.09}{\xR-0.09}{-3.4}{speichern(a)}
\nachricht{\xR+0.09}{\xB-0.09}{-3.95}{transaktion()}
\rueckgabeL{\xB-0.09}{\xR+0.09}{-4.45}{INSERT bzw. UPDATE, Commit}
\nachricht[antwort]{\xR-0.09}{\xS+0.09}{-4.95}{a (mit id)}
\nachricht[antwort]{\xS-0.09}{\xD+0.09}{-5.45}{a}
\nachricht{\xD+0.09}{\xZ-0.09}{-6.0}{daten\_geaendert.emit()}
\nachricht{\xZ+0.09}{\xH-0.09}{-6.55}{\_geaendert()}
\selbst{\xH}{-6.85}{[sichtbar] aktualisieren()}
\aufrufR{\xH+0.09}{\xA-0.09}{-7.75}{wochenbewertung(kw)}
\aufrufR{\xA+0.09}{\xR-0.09}{-8.3}{im\_zeitraum(montag, sonntag)}
\rueckgabeL{\xR-0.09}{\xA+0.09}{-8.8}{aktivitaeten : list[Aktivitaet]}
\aufrufR{\xA+0.09}{\xW-0.09}{-9.5}{bewerte\_woche(\dots)}
\rueckgabeL{\xW-0.09}{\xA+0.09}{-10.0}{Wochenbewertung}
\rueckgabeL{\xA-0.09}{\xH+0.09}{-10.35}{Wochenbewertung}
\nachricht{\xH+0.09}{\xA-0.09}{-11.1}{hinweise(kw)}
\aufrufR{\xA+0.09}{\xG-0.09}{-11.45}{pruefe\_regeln(k)}
\node[notiz, anchor=north west, text width=3.1cm, execute at begin node=\hyphenpenalty 10000\relax, font=\fontsize{6.8pt}{7.8pt}\selectfont]
  at (\xS-0.2,-10.75) {hinweise() lädt zuvor ebenfalls Aktivitäten (drei Wochen) und bewertet sie;
  diese Aufrufe sind hier weggelassen.};
\nachricht[antwort]{\xG-0.09}{\xA+0.09}{-11.95}{list[Hinweis]}
\nachricht[antwort]{\xA-0.09}{\xH+0.09}{-12.45}{list[Hinweis]}
\selbst{\xH}{-12.75}{Balken und Hinweise anzeigen}
\selbst{\xD}{-13.42}{accept()}

% sonst: Eingaben ungültig
\rueckgabeL{\xS-0.09}{\xD+0.09}{-14.6}{raise ValidierungsFehler}
\selbst{\xD}{-15.0}{\_fehler\_anzeigen(e.fehler)}
\nachricht[antwort]{\xD-0.09}{\xN}{-15.75}{Meldung am Feld}

\end{tikzpicture}
\end{document}
